\section{Kết quả thực nghiệm và thảo luận}
\label{sec:thuc_nghiem}

\subsection{Thiết lập thực nghiệm}
Mô phỏng được xây dựng trên bộ dữ liệu thực tế nhằm đánh giá độ hội tụ và tính ổn định của giải thuật. Kết quả so sánh giữa phương pháp đề xuất và các phương pháp chuẩn (Baseline) được trình bày chi tiết.

\begin{table}[H]
  \centering
  \caption{So sánh hiệu năng giữa giải thuật đề xuất và các Baseline}
  \label{tab:results}
  \begin{tabular}{lcccc}
    \toprule
    \textbf{Phương pháp} & \textbf{Độ chính xác (\%)} & \textbf{Thời gian (s)} & \textbf{Năng lượng (J)} & \textbf{Cải thiện} \\
    \midrule
    Baseline A & 82.4 & 145.2 & 320.5 & Gốc \\
    Baseline B & 86.1 & 112.8 & 275.0 & +4.5\% \\
    \textbf{Đề xuất (Ours)} & \textbf{91.8} & \textbf{78.4} & \textbf{189.2} & \textbf{+11.4\%} \\
    \bottomrule
  \end{tabular}
\end{table}

\subsection{Thảo luận kết quả}
Từ Bảng \ref{tab:results}, ta nhận thấy phương pháp đề xuất giảm thời gian xử lý tới $46.0\%$ so với Baseline A trong khi vẫn nâng cao độ chính xác lên $91.8\%$. Điều này minh chứng tính đúng đắn của giả thuyết nghiên cứu đã đặt ra ban đầu.
